% ---
% titre: Loto du crible d'Eratosthene
% theme: NO
% type: JEU
% chapitre: Nombres naturels et decimaux
% sous_chapitre: Nombres premiers
% niveau: 9
% difficulte: 1
% corrige: oui
% tags: c-nombres-premiers, c-multiples, f-commun, p-decouverte, p-institutionnalisation
% source: creation originale
% ---

% ============================================================
% FICHE ENSEIGNANT-E
% ============================================================

\section*{Loto du crible d'\'Eratosth\`ene --- Fiche enseignant\textperiodcentered e}

\begin{tcolorbox}[colback=white, colframe=themeColor, title=Objectif, fonttitle=\bfseries]
Faire d\'ecouvrir les nombres premiers et le principe du crible d'\'Eratosth\`ene
\`a travers un jeu de type loto, sans annoncer la r\`egle \`a l'avance : ce sont
les \'el\`eves qui, en comparant leurs grilles \`a la fin, d\'ecouvrent que le
"tirage al\'eatoire" ne l'\'etait pas.
\end{tcolorbox}

\begin{tcolorbox}[colback=white, colframe=themeColor, title=Mat\'eriel, fonttitle=\bfseries]
\begin{itemize}
  \item Une grille par \'el\`eve, \`a choisir parmi les variantes A, B, C ci-dessous
  (m\^emes nombres de 1 \`a 100, mise en page diff\'erente) --- \textbf{r\'epartir
  les variantes dans la classe de mani\`ere m\'elang\'ee}, de sorte que deux
  voisin\textperiodcentered es directs n'aient si possible pas la m\^eme variante.
  \item Un sac ou une bo\^ite opaque contenant des papiers num\'erot\'es (voir
  \emph{ordre des tirages} ci-dessous) --- ou simplement une liste que
  l'enseignant\textperiodcentered e annonce \`a voix haute sans montrer de sac.
  \item Deux crayons de couleur par \'el\`eve (une couleur pour entourer un
  survivant, une couleur pour rayer un multiple).
\end{itemize}
\end{tcolorbox}

\begin{tcolorbox}[colback=white, colframe=themeColor, title=R\`egle annonc\'ee aux \'el\`eves, fonttitle=\bfseries]
\begin{enumerate}
  \item Je tire un nombre.
  \item S'il n'est pas encore ray\'e sur votre grille, entourez-le en vert :
  c'est un "survivant".
  \item Rayez ensuite en rouge tous ses multiples sur la grille (jusqu'\`a 100),
  s'ils ne le sont pas d\'ej\`a.
  \item Si le nombre tir\'e est d\'ej\`a ray\'e, il ne se passe rien --- attendez
  le tirage suivant.
\end{enumerate}
Ne pas pr\'eciser que seuls des nombres premiers seront tir\'es.
\end{tcolorbox}

\begin{tcolorbox}[colback=white, colframe=themeColor, title=Ordre exact des tirages, fonttitle=\bfseries]
\begin{tabularx}{\linewidth}{c c X}
\textbf{\#} & \textbf{Tirage} & \textbf{Ce qu'il se passe} \\ \hline
1 & 2 & Premier survivant. Rayer tous les multiples de 2 (50 nombres). \\
2 & 3 & Survivant. Rayer les multiples de 3 non d\'ej\`a ray\'es. \\
3 & 9 \emph{(leurre)} & D\'ej\`a ray\'e par le tirage pr\'ec\'edent --- rien de nouveau. Bon moment pour demander "pourquoi personne ne bouge ?" \\
4 & 5 & Survivant. Rayer les multiples de 5 restants. \\
5 & 7 & Survivant. Rayer les multiples de 7 restants (peu de nouvelles cases). \\
6 & 11 \emph{(question, pas un vrai tirage)} & Demander : "si je tire 11, est-ce que \c{c}a va encore rayer des cases ?" Faire v\'erifier : $11\times2, 11\times3,\dots$ sont d\'ej\`a tous ray\'es en dessous de 100. \\
\end{tabularx}
\end{tcolorbox}

\begin{tcolorbox}[colback=white, colframe=themeColor, title=Moment de bascule --- crit\`ere d'arr\^et, fonttitle=\bfseries]
Faire chercher pourquoi 11 (et tout nombre premier sup\'erieur) ne peut plus
rien \'eliminer de nouveau en dessous de 100 : tout multiple de 11 inf\'erieur
\`a 100 s'\'ecrit $11 \times k$ avec $k \le 9$, donc $k$ a d\'ej\`a \'et\'e
trait\'e (2, 3, 5 ou 7 le divise, ou $k=1$ redonne 11 lui-m\^eme). Amener
ensuite l'id\'ee g\'en\'erale : on peut s'arr\^eter d\`es que le nombre tir\'e
d\'epasse $\sqrt{100}=10$.
\end{tcolorbox}

\begin{tcolorbox}[colback=white, colframe=themeColor, title=R\'ev\'elation finale, fonttitle=\bfseries]
"Comparez la liste de vos nombres entour\'es en vert avec votre
voisin\textperiodcentered e." M\^eme si les grilles ont des
formes diff\'erentes, la liste des survivants est identique pour toute la
classe : 2, 3, 5, 7, 11, 13, 17, 19, 23, 29, 31, 37, 41, 43, 47, 53, 59, 61,
67, 71, 73, 79, 83, 89, 97. Ce n'\'etait pas un vrai loto al\'eatoire --- ce
sont les nombres premiers, d\'efinis par une propri\'et\'e fixe et non par le
hasard du tirage. Institutionnaliser ensuite la d\'efinition de nombre premier
et le principe du crible d'\'Eratosth\`ene.
\end{tcolorbox}

\begin{tcolorbox}[colback=white, colframe=themeColor, title=Points de vigilance, fonttitle=\bfseries]
\begin{itemize}
  \item Clarifier avant de commencer que \textbf{1 n'est ni premier ni
  compos\'e} : il reste entour\'e en gris ou simplement ignor\'e, pour \'eviter
  les grilles bloqu\'ees sur ce cas particulier.
  \item Garder un rythme soutenu entre les tirages (1--2 minutes max) pour ne
  pas perdre l'effet ludique pendant que les \'el\`eves plus lent\textperiodcentered
  es terminent de rayer.
  \item R\'epartir vraiment les 3 variantes de grille dans la classe, sinon des
  \'el\`eves voisin\textperiodcentered es risquent de rep\'erer la co\"incidence
  trop t\^ot en regardant la feuille du voisin.
\end{itemize}
\end{tcolorbox}

\clearpage

% ============================================================
% GRILLE ELEVE --- VARIANTE A (grille classique 10x10)
% ============================================================
\header{Loto du crible d'Ératosthène }

Le nombre annoncé a survécu jusqu'ici ? Entoure-le en vert, c'est un survivant ! Puis élimine (raye en rouge) tous ses multiples.

\vfill
\begin{center}
\begin{tikzpicture}[x=1.5cm,y=1.5cm]
\foreach \n in {1,...,100}{
  \pgfmathtruncatemacro{\col}{mod(\n-1,10)}
  \pgfmathtruncatemacro{\row}{9-int((\n-1)/10)}
  \draw[draw=themeColor, thin] (\col,\row) rectangle ++(1,1);
  \node[font=\small] at (\col+0.5,\row+0.5) {\n};
}
\end{tikzpicture}
\end{center}
\vfill

\newpage
% ============================================================
% GRILLE ELEVE --- VARIANTE B (meme grille, presentee en losange)
% ============================================================
\header{Loto du crible d'Ératosthène }

Le nombre annoncé a survécu jusqu'ici ? Entoure-le en vert, c'est un survivant ! Puis élimine (raye en rouge) tous ses multiples.


\vfill
\begin{center}
\begin{tikzpicture}[x=1.2cm,y=1.2cm]
\begin{scope}[rotate=45]
\foreach \n in {1,...,100}{
  \pgfmathtruncatemacro{\col}{mod(\n-1,10)}
  \pgfmathtruncatemacro{\row}{9-int((\n-1)/10)}
  \draw[draw=themeColor, thin] (\col,\row) rectangle ++(1,1);
  \node[font=\small] at (\col+0.5,\row+0.5) {\n};
}
\end{scope}
\end{tikzpicture}
\end{center}
\vfill

\newpage
% ============================================================
% GRILLE ELEVE --- VARIANTE C (meme grille, format 5x20 en cercles)
% ============================================================
\header{Loto du crible d'Ératosthène }

Le nombre annoncé a survécu jusqu'ici ? Entoure-le en vert, c'est un survivant ! Puis élimine (raye en rouge) tous ses multiples.

\vfill
\begin{center}
\begin{tikzpicture}[x=3cm,y=1cm]
\foreach \n in {1,...,100}{
  \pgfmathtruncatemacro{\col}{mod(\n-1,5)}
  \pgfmathtruncatemacro{\row}{19-int((\n-1)/5)}
  \draw[draw=themeColor, thin] (\col,\row) circle (0.42);
  \node[font=\small] at (\col,\row) {\n};
}
\end{tikzpicture}
\end{center}
\vfill
